% ============================================================
% tab/tabla_comparativa.tex
%
% COLUMNAS OBLIGATORIAS: referencia, año, técnica, dataset o
% entorno, métrica principal, resultado reportado, limitación.
%
% AJUSTE DE ANCHOS
%   Referencia y Año: anchos fijos, porque su contenido es corto
%   y predecible. Las cinco restantes son columnas X con peso,
%   de modo que el espacio sobrante se reparte en proporción y
%   la tabla se adapta si cambia el ancho de página.
%   Los pesos suman 5, que es el número de columnas W.
%
%   NO usar 'l' en la primera columna: no permite salto de línea
%   y cualquier marcador largo desborda la tabla.
% ============================================================
\begin{table*}[!t]
\centering
\caption{Resumen comparativo de trabajos relacionados.}
\label{tab:comparativa}
\begin{threeparttable}
\footnotesize
\setlength{\tabcolsep}{4pt}
\begin{tabularx}{\textwidth}{@{}
    P{1.3cm}                        % Referencia
    >{\centering\arraybackslash}p{0.7cm} % Año
    W{0.95}                         % Técnica
    W{1.00}                         % Dataset / entorno
    W{0.85}                         % Métrica principal
    W{0.75}                         % Resultado reportado
    W{1.45}                         % Limitación identificada
  @{}}
\toprule
\textbf{Referencia} & \textbf{Año} & \textbf{Técnica} &
\textbf{Dataset / entorno} & \textbf{Métrica principal} &
\textbf{Resultado reportado} & \textbf{Limitación identificada}\\
\midrule

\cite{hodge_winprediction_2021} & 2021 &
Regresión logística, bosques aleatorios y potenciación del gradiente &
DotA 2: 5\,744 repeticiones mixtas y 186 profesionales &
Exactitud &
\pct &
Los conjuntos de la literatura no son comparables entre sí; el flujo en
vivo provee un subconjunto de las características de las repeticiones
\\
\addlinespace[3pt]

\cite{xenopoulos_winprob_2022} & 2022 &
Máquinas de potenciación explicables (modelo aditivo generalizado) &
CSGO: tres entornos de habilidad &
Pérdida logarítmica y error de calibración esperado &
\pct &
Un solo juego; el modelo interpretable podría ser superado; no es posible
identificar a los jugadores que incurren en trampa
\\
\addlinespace[3pt]

\cite{saravanan_metadiscovery_2024} & 2024 &
Agente de refuerzo y agente heurístico, con constructor de equipos &
Pok\'emon Showdown: cuatro escenarios de prohibición &
Solapamiento, distancia de edición y correlación de rangos &
\pct &
La correlación de rangos resulta débil y no significativa: el conjunto
descubierto difiere del real pese al solapamiento
\\
\addlinespace[3pt]

\pct & \pct & \pct & \pct & \pct & \pct & \pct \\
\addlinespace[3pt]

\pct & \pct & \pct & \pct & \pct & \pct & \pct \\
\addlinespace[3pt]

\pct & \pct & \pct & \pct & \pct & \pct & \pct \\
\addlinespace[3pt]

\pct & \pct & \pct & \pct & \pct & \pct & \pct \\

\bottomrule
\end{tabularx}
\begin{tablenotes}[flushleft]
\footnotesize
\item \textit{Nota.} Los conjuntos de datos empleados difieren en tamaño,
composición, juego y versión, por lo que la comparación directa de los
resultados reportados requiere cautela. Las referencias de carácter
metodológico incluidas en la revisión no figuran en esta tabla, dado que
no resuelven una tarea de predicción sobre un conjunto de datos
comparable.
\end{tablenotes}
\end{threeparttable}
\end{table*}